\section{Multivariate OTI numbers and their real matrix representation}
\label{sec:multivariate}
Sections~\ref{sec:oti} and~\ref{sec:cr} use one parameter and one nilpotent unit.
This appendix extends the construction to $d$ independent parameters and
explains how one evaluation can carry pure and mixed derivatives. The algebra
and matrix representation follow the multivariate OTI construction
of~\cite{aristizabal2020}; the basis ordering is stated explicitly below.

\subsection{Algebra, seeding, and mixed derivatives}
\label{sec:multi-algebra}
Introduce commuting units $\varepsilon_1,\ldots,\varepsilon_d$ and retain
monomials of \emph{total degree} at most $p$. For a multi-index
$\alpha=(\alpha_1,\ldots,\alpha_d)$ of nonnegative integers, define
\[
 |\alpha|=\alpha_1+\cdots+\alpha_d,\qquad
 \alpha!=\alpha_1!\cdots\alpha_d!,\qquad
 \boldsymbol\varepsilon^\alpha=
 \varepsilon_1^{\alpha_1}\cdots\varepsilon_d^{\alpha_d}.
\]
A multivariate OTI number is
\begin{equation}
\label{eq:multi-number}
 a=\sum_{|\alpha|\leq p}a_\alpha\boldsymbol\varepsilon^\alpha,
 \qquad a_\alpha\in\mathbb R,
 \qquad \boldsymbol\varepsilon^\alpha=0\ \text{if }|\alpha|>p.
\end{equation}
For a second number $b=\sum_{|\alpha|\leq p}b_\alpha\boldsymbol\varepsilon^\alpha$,
addition combines equal monomials, while multiplication adds their exponents
and discards products whose total degree exceeds $p$:
\begin{equation}
\label{eq:multi-arithmetic}
 a+b=\sum_{|\alpha|\leq p}(a_\alpha+b_\alpha)
           \boldsymbol\varepsilon^\alpha,\qquad
 ab=\sum_{|\gamma|\leq p}
       \left(\sum_{\alpha+\beta=\gamma}a_\alpha b_\beta\right)
       \boldsymbol\varepsilon^\gamma.
\end{equation}
For two variables at order two, the retained monomials are
$1,\varepsilon_1,\varepsilon_2,\varepsilon_1^2,
\varepsilon_1\varepsilon_2,\varepsilon_2^2$.
The mixed term $\varepsilon_1\varepsilon_2$ is retained, but
$\varepsilon_1^2\varepsilon_2$ and $\varepsilon_1\varepsilon_2^2$ vanish.
Requiring only $\varepsilon_i^{p+1}=0$ separately would not impose this
\emph{total-degree} truncation: mixed monomials above degree $p$ must also vanish.

For an analytic $f(x_1,\ldots,x_d)$, seed each parameter independently as
$x_i^*=x_{i,0}+\varepsilon_i$. Its multivariate Taylor expansion is
\begin{equation}
\label{eq:multi-taylor}
 f(\boldsymbol x_0+\boldsymbol\varepsilon)
 =\sum_{|\alpha|\leq p}
   \frac{\partial^\alpha f(\boldsymbol x_0)}{\alpha!}
   \boldsymbol\varepsilon^\alpha,\qquad
 \partial^\alpha f=
 \frac{\partial^{|\alpha|}f}
 {\partial x_1^{\alpha_1}\cdots\partial x_d^{\alpha_d}}.
\end{equation}
Extending the coefficient-extraction notation of Equation~\eqref{eq:oti-extraction},
\begin{equation}
\label{eq:multi-extraction}
 \boxed{\partial^\alpha f(\boldsymbol x_0)=
 \alpha!\,\operatorname{Im}_{\boldsymbol\varepsilon^\alpha}
 \bigl(f(\boldsymbol x_0+\boldsymbol\varepsilon)\bigr)}.
\end{equation}
The scale is the product of factorials $\alpha!$, not the factorial of the
total degree $|\alpha|$. For example, the coefficient of $\varepsilon_1^2$
is $f_{x_1x_1}/2!$, whereas the coefficient of
$\varepsilon_1\varepsilon_2$ is $f_{x_1x_2}/(1!1!)=f_{x_1x_2}$.
Seeding both inputs with the same unit would instead give a directional
derivative and would not keep their partial derivatives separate.

\newpage
\subsection{Building the matrix algebra from its generators}
\label{sec:multi-matrices}
The OTI algebra of Section~\ref{sec:multi-algebra} has $d$ independent
commuting generators $\varepsilon_1,\ldots,\varepsilon_d$ and truncation
order $p$. The generators produce all higher-order monomials by
multiplication. The full basis consists of the constant monomial $1$
and all products of these generators through total degree $p$.
To build a matrix representation, it is therefore sufficient to construct
one matrix $N_i$ representing each generator $\varepsilon_i$.
The matrix for a higher-order monomial is then the corresponding product:
\begin{equation}
 \label{eq:multi-monomial-matrices}
 1\longmapsto I_q,\qquad
 \varepsilon_i\longmapsto N_i,\qquad
 \varepsilon_1^{\alpha_1}\cdots\varepsilon_d^{\alpha_d}
 \longmapsto N_1^{\alpha_1}\cdots N_d^{\alpha_d},
\end{equation}
where $q=\binom{d+p}{p}$ is the number of retained basis monomials.
These matrices must reproduce the algebra's defining relations: the $N_i$
commute, and every product of more than $p$ generating factors is zero.
The construction below enforces both properties.

\paragraph{Construct each generator matrix by its action on the basis.}
Choose an ordering of the retained monomials, beginning with $1$, and let
$[a]$ denote the coordinate vector representing an OTI number $a$ in this
order. The number itself remains the algebraic object in
Equation~\eqref{eq:multi-number}. We define $N_i$ to perform multiplication
by $\varepsilon_i$:
\begin{equation}
 \label{eq:multi-shift}
 \boxed{N_i[a]=[\varepsilon_i a].}
\end{equation}
To determine column $j$, multiply the $j$th basis monomial by
$\varepsilon_i$ and express the result in the chosen coordinates.
If the resulting total degree exceeds $p$, that column is zero.
Thus every column follows directly from multiplication in the OTI algebra.

\paragraph{Represent a general number by substitution.}
Once the generator matrices are known, keep each real coefficient and
replace its monomial by the matrix product in
Equation~\eqref{eq:multi-monomial-matrices}:
\begin{equation}
\label{eq:multi-cr}
 \boxed{C(a)=\sum_{|\alpha|\leq p}a_\alpha
 N_1^{\alpha_1}\cdots N_d^{\alpha_d}
 =\sum_{|\alpha|\leq p}a_\alpha N^\alpha},
 \qquad N^{\boldsymbol0}=I_q.
\end{equation}
Here $N^\alpha$ abbreviates the product of generator matrices. Since the
generators commute, these products add multi-indices,
$N^\alpha N^\beta=N^{\alpha+\beta}$, with products above total degree $p$
equal to zero, exactly as for the monomials $\boldsymbol\varepsilon^\alpha$.
Consequently, for any two OTI numbers $a$ and $b$ in the same algebra,
\begin{equation}
 \label{eq:multi-multiplication-action}
 C(a+b)=C(a)+C(b),\qquad C(ab)=C(a)C(b).
\end{equation}
The first column of $C(a)$ holds the coefficients $a_\alpha$ in the chosen
monomial order. The independently seeded input
$x_i^*=x_{i,0}+\varepsilon_i$ is represented by $X_i=x_{i,0}I_q+N_i$.
As in Section~\ref{sec:cr}, analytic elementary functions are evaluated
as matrix functions of these representations.

\newpage
\paragraph{Example: two generators and truncation order two.}
For $d=2$, $p=2$, choose the ordered basis
\[
 \mathcal B=(1,\varepsilon_1,\varepsilon_2,\varepsilon_1^2,
 \varepsilon_1\varepsilon_2,\varepsilon_2^2).
\]
Multiplying each entry of this tuple by a generator gives
\begin{align}
 \varepsilon_1\mathcal B
 &=(\varepsilon_1,\varepsilon_1^2,\varepsilon_1\varepsilon_2,0,0,0),
 \notag\\
 \varepsilon_2\mathcal B
 &=(\varepsilon_2,\varepsilon_1\varepsilon_2,\varepsilon_2^2,0,0,0).
 \label{eq:multi-basis-action}
\end{align}
The last three entries vanish because their total degree would be three.
Each output supplies one column of the corresponding matrix, expressed
in the coordinates of $\mathcal B$. Consequently,
\[
 N_1=\begin{bmatrix}
 0&0&0&0&0&0\\1&0&0&0&0&0\\0&0&0&0&0&0\\
 0&1&0&0&0&0\\0&0&1&0&0&0\\0&0&0&0&0&0
 \end{bmatrix},\qquad
 N_2=\begin{bmatrix}
 0&0&0&0&0&0\\0&0&0&0&0&0\\1&0&0&0&0&0\\
 0&0&0&0&0&0\\0&1&0&0&0&0\\0&0&1&0&0&0
 \end{bmatrix}.
\]
For example, $\varepsilon_1\cdot\varepsilon_2=\varepsilon_1\varepsilon_2$:
the third input monomial maps to the fifth basis monomial, placing a one
in row five, column three of $N_1$. The higher-order monomials are
represented by $N_1^2$, $N_1N_2=N_2N_1$, and $N_2^2$.
Thus there are two generating matrices and six basis matrices,
\[
 I_6,\quad N_1,\quad N_2,\quad N_1^2,\quad N_1N_2,\quad N_2^2.
\]
For a general OTI number
\[
 a=a_{00}+a_{10}\varepsilon_1+a_{01}\varepsilon_2+
 a_{20}\varepsilon_1^2+a_{11}\varepsilon_1\varepsilon_2+a_{02}\varepsilon_2^2,
\]
Equation~\eqref{eq:multi-cr} becomes
\begin{equation}
 \label{eq:multi-cr-generators}
 \boxed{C(a)=a_{00}I_6+a_{10}N_1+a_{01}N_2
 +a_{20}N_1^2+a_{11}N_1N_2+a_{02}N_2^2.}
\end{equation}
Expanding this linear combination gives
\begin{equation}
\label{eq:multi-cr-explicit}
 C(a)=\begin{bmatrix}
 a_{00}&0&0&0&0&0\\
 a_{10}&a_{00}&0&0&0&0\\
 a_{01}&0&a_{00}&0&0&0\\
 a_{20}&a_{10}&0&a_{00}&0&0\\
 a_{11}&a_{01}&a_{10}&0&a_{00}&0\\
 a_{02}&0&a_{01}&0&0&a_{00}
 \end{bmatrix}.
\end{equation}
The first column agrees with $[a]$, as required. Every other column records
multiplication of $a$ by the corresponding basis monomial. These matrices
form a six-dimensional subalgebra of the real $6\times6$ matrices.

\newpage
\subsection{Worked example with pure and mixed derivatives}
\label{sec:multi-example}
Extend the scalar example of Section~\ref{sec:example-oti} to
$f(x,y)=\exp(x^2+y)$ at $(x_0,y_0)=(1,0)$, retaining total degree two.
The independently seeded inputs and the argument of the exponential are
\[
 x^*=1+\varepsilon_1,\qquad y^*=\varepsilon_2,\qquad
 (x^*)^2+y^*=1+2\varepsilon_1+\varepsilon_1^2+\varepsilon_2.
\]
The Taylor expansion of the exponential about its real part 1 gives
\begin{align*}
 f(x^*,y^*)&=\mathrm e\left[1+(2\varepsilon_1+\varepsilon_1^2+\varepsilon_2)
 +\frac{(2\varepsilon_1+\varepsilon_1^2+\varepsilon_2)^2}{2}\right],\\
 (2\varepsilon_1+\varepsilon_1^2+\varepsilon_2)^2
 &=4\varepsilon_1^2+4\varepsilon_1\varepsilon_2+\varepsilon_2^2.
\end{align*}
All omitted products have total degree at least three. Therefore
\begin{equation}
\label{eq:multi-example-oti}
 f(x^*,y^*)=\mathrm e\left(1+2\varepsilon_1+\varepsilon_2+
 3\varepsilon_1^2+2\varepsilon_1\varepsilon_2+
 \tfrac12\varepsilon_2^2\right).
\end{equation}
For the real matrix evaluation, use $X=I_6+N_1$ and $Y=N_2$. The identical
sequence of operations is a matrix product followed by a matrix exponential:
\begin{align}
 V&=X^2+Y=I_6+2N_1+N_1^2+N_2,\notag\\
 W&=\exp(V)=\mathrm e\left(I_6+2N_1+N_2+3N_1^2+
 2N_1N_2+\tfrac12N_2^2\right).
 \label{eq:multi-example-matrix}
\end{align}
The second equality follows by expanding about $I_6$; products of three
nilpotent factors vanish. The first column is
\[
 W\boldsymbol e_{\boldsymbol0}
 =\begin{bmatrix}\mathrm e&2\mathrm e&\mathrm e&3\mathrm e&2\mathrm e&\mathrm e/2\end{bmatrix}^{T}.
\]
Reading the entries in the declared monomial order and applying
Equation~\eqref{eq:multi-extraction} yields
\begin{equation}
\label{eq:multi-example-derivatives}
 f(1,0)=\mathrm e,\qquad
 \nabla f(1,0)=\begin{bmatrix}2\mathrm e\\\mathrm e\end{bmatrix},\qquad
 \nabla^2 f(1,0)=\begin{bmatrix}6\mathrm e&2\mathrm e\\2\mathrm e&\mathrm e\end{bmatrix}.
\end{equation}
In particular, the mixed derivative is the fifth entry without an extra
factor of two, while the fourth and sixth entries are multiplied by $2!$.
Direct differentiation of $\exp(x^2+y)$ confirms all these values.

\newpage
\subsection{General matrices with multivariate OTI entries}
\label{sec:multi-matrix-arguments}
The Kronecker construction of Equation~\eqref{eq:matrix-cr-kronecker} extends
by replacing powers of one shift matrix with products of the $N_i$:
\begin{equation}
\label{eq:multi-block-cr}
 M=\sum_{|\alpha|\leq p}M_\alpha\boldsymbol\varepsilon^\alpha
 \quad\longmapsto\quad
 \boxed{\mathcal C(M)=\sum_{|\alpha|\leq p}N^\alpha\otimes M_\alpha}.
\end{equation}
If $M_\alpha\in\mathbb R^{m\times m}$, the representation is a single real
$qm\times qm$ matrix. The coefficient matrices may be arbitrary and need not
commute. Their order is preserved by
$(N^\alpha\otimes M_\alpha)(N^\beta\otimes L_\beta)
 =N^{\alpha+\beta}\otimes(M_\alpha L_\beta)$,
with products above total degree $p$ equal to zero. Consequently, the
representation preserves matrix multiplication and the matrix exponential.

For example, let $M(x,y)=xA+yB$, with fixed real $m\times m$ matrices $A,B$,
and let $D=x_0A+y_0B$. Seeding both parameters at total degree two gives
\begin{equation}
\label{eq:multi-affine-block}
 \mathcal M=I_6\otimes D+N_1\otimes A+N_2\otimes B
 =\begin{bmatrix}
 D&0&0&0&0&0\\
 A&D&0&0&0&0\\
 B&0&D&0&0&0\\
 0&A&0&D&0&0\\
 0&B&A&0&D&0\\
 0&0&B&0&0&D
 \end{bmatrix}.
\end{equation}
Every displayed entry is an $m\times m$ block. For
$F(x,y)=\exp(xA+yB)$, the first block column of $\exp(\mathcal M)$ is
\begin{equation}
\label{eq:multi-matrix-extraction}
 \left.\begin{bmatrix}
 F\\F_x\\F_y\\F_{xx}/2!\\F_{xy}\\F_{yy}/2!
 \end{bmatrix}\right|_{(x_0,y_0)}.
\end{equation}
No assumption that $AB=BA$ is required; in general the derivatives cannot be
simplified to products such as $A\exp(D)$ or $AB\exp(D)$.

\paragraph{Verification and scope.}
The companion script below verifies the two-variable multiplication matrices,
the scalar example in Equation~\eqref{eq:multi-example-derivatives}, and the
block construction with a noncommuting pair $A,B$:
\begin{verbatim}
julia --project=. docs/verify_multivariate_example.jl
\end{verbatim}
The thermal sensitivity entry points remain single-parameter implementations.
This appendix supplies the multivariate construction and a separate numerical
check. For $d$ parameters at total order $p$, $q=\binom{d+p}{p}$ grows with both
$d$ and $p$, so the dense real representation becomes substantially larger.
It should not be replaced by an unmodified tensor product of univariate
representations: that would retain terms up to degree $p$ in each variable,
rather than the total-degree truncation specified in Section~\ref{sec:multi-algebra}.
